\begin{abstract}
A text embedding teacher stores knowledge in the geometry it induces over an
entire corpus, but supervising all pairwise relations is infeasible. Standard
relational knowledge distillation (KD) instead forms its supervision support
from examples that happen to co-occur in a mini-batch, so relation exposure is
independent of semantic relatedness: a related pair and an unrelated corpus
pair are equally likely to be supervised. We argue that a
corpus-level relational objective should separate two decisions that mini-batch
KD conflates---which relations are \emph{eligible} for supervision, and how
eligible relations are \emph{realized} efficiently---and that preserving global
geometry depends on both alignment, reproducing the supervised relations
accurately, and coverage, supervising most of the relational mass the teacher
considers important. We introduce \textbf{Global Geometry Preservation
Knowledge Distillation (GGPKD)}, which freezes the base student's $k$NN
shortlists, asks the teacher to value the selected relations, and uses
mini-batches only to realize them. GGPKD encodes the corpus once with the
teacher and base student, then forms teacher-valued transition targets on the
student-selected columns. The deduplicated candidate pool of a batch is encoded
once and reused twice more: to calibrate similarity levels across the batch,
and to supervise additional graph-selected rows without any extra student
encoding. Independent coefficients weight the ambient and graph relational
groups directly, alongside the graph width and auxiliary-row coefficient.
Across nine standard sentence-embedding benchmarks and
three teacher--student settings, GGPKD obtains the best mean score in every
setting, with the largest and most consistent gains on semantic textual
similarity. GGPKD requires no task labels and no online teacher inference, and
only the student encoder is used at deployment.
\end{abstract}
